\section{Methodological Framework: Bounding the Relaxation Gap}
\label{sec:framework}

The continuous prompt space becomes useful only if relaxed optimization can be
related to the quality of the projected discrete attack. We make that
connection through a Lipschitz regularity assumption on the target and draft
distributions.

Let $p_i^U$ and $q_i^U$ denote the target and draft next-token distributions at
speculative step $i$ under relaxed prefix $U$. Consider the weighted
disagreement objective
\[
C(U) = \sum_{i=1}^{\gamma} w_i \tv(p_i^U, q_i^U),
\]
where $w_i = \gamma - i + 1$ linearly emphasizes early speculative steps. The
corresponding discrete objective is
\[
C_d(z) := C(E(z)).
\]
To recover a deployable attack, define the token-wise nearest-neighbor
projection $\Pi(U) := (\Pi_1(U), \dots, \Pi_m(U))$ by
\[
\Pi_t(U) = \arg\min_{v \in \mathcal{V}} \|e_t(U) - E_v\|_2,
\]
and define the projection distortion
\[
\rho(U) := \sum_{t=1}^{m} \|e_t(U) - E_{\Pi_t(U)}\|_2.
\]

\begin{assumption}[TV-Lipschitzness]
\label{assump:tv-lipschitz}
For each speculative step $i$, there exist constants $L_i^p, L_i^q$ such that
for any two relaxed prefixes $U, V$,
\[
\tv(p_i^U, p_i^V) \le L_i^p \|U - V\|_{2,1},
\qquad
\tv(q_i^U, q_i^V) \le L_i^q \|U - V\|_{2,1},
\]
where
\[
\|U - V\|_{2,1} := \sum_{t=1}^{m} \|e_t(U) - e_t(V)\|_2.
\]
Let
\[
L_C := \sum_{i=1}^{\gamma} w_i (L_i^p + L_i^q)
\]
denote the aggregate Lipschitz constant of the weighted objective.
\end{assumption}

Under this assumption, the discrete quality of a projected relaxed solution can
be bounded explicitly.

\begin{theorem}[Relaxation Gap for Early Acceptance Collapse]
\label{thm:relax-gap}
Let $U^*$ be an $\varepsilon$-optimal solution of the relaxed problem such that
\[
C(U^*) \ge \sup_U C(U) - \varepsilon.
\]
Let
\[
z^* \in \arg\max_{z \in \mathcal{V}^m} C_d(z)
\]
be the optimal discrete prefix. Then
\[
C_d(\Pi(U^*)) \ge C_d(z^*) - \varepsilon - L_C \rho(U^*).
\]
\end{theorem}

\begin{proof}
By Assumption~\ref{assump:tv-lipschitz}, the objective $C$ is $L_C$-Lipschitz
under $\|\cdot\|_{2,1}$. Let $\hat{z} = \Pi(U^*)$. Since the continuous
embeddings of $\hat{z}$ are within $\rho(U^*)$ of $U^*$, we have
\[
C_d(\hat{z}) = C(E(\hat{z})) \ge C(U^*) - L_C \rho(U^*).
\]
Because the optimal discrete prefix is feasible in the relaxed space,
\[
C(U^*) \ge C_d(z^*) - \varepsilon.
\]
Combining the two inequalities yields the claim.
\end{proof}

The theorem isolates two quantities that determine whether relaxed search
survives projection: the solver suboptimality $\varepsilon$ and the
projection distortion $\rho(U^*)$.

We can also state an exact-recovery condition. Define the discrete attack
margin
\[
\Gamma := C_d(z^*) - \max_{z \neq z^*} C_d(z).
\]

\begin{corollary}[Exact Recovery under Discrete Margin]
\label{cor:exact-recovery}
Assume the optimal discrete adversarial prefix $z^*$ is unique
($\Gamma > 0$). If
\[
\varepsilon + L_C \rho(U^*) < \Gamma,
\]
then
\[
\Pi(U^*) = z^*.
\]
\end{corollary}

These results directly shape the optimizer. Maximizing the relaxed objective
alone is insufficient; the solver must also control the distance from the
relaxed solution to the token manifold. This motivates a penalized objective of
the form
\[
\max_U \; C(U) - \lambda \rho(U) - \tau \sum_{t=1}^{m} H(u_t),
\]
where $\lambda \rho(U)$ is a projection-gap control term and the entropy term
sharpens the simplex vectors toward projectable token configurations. ADSD is
the first attack algorithm in our setting designed explicitly around this
relaxation-gap logic.
